\documentclass[10pt,a4paper]{article}
\usepackage[margin=1.7cm]{geometry}
\usepackage[T1]{fontenc}
\usepackage[utf8]{inputenc}
\usepackage{lmodern}
\usepackage{graphicx}
\usepackage{enumitem}
\usepackage{titlesec}
\usepackage[hidelinks,unicode]{hyperref}
\input{glyphtounicode}
\pdfgentounicode=1
\pagestyle{empty}
\setlength{\parindent}{0pt}
\setlength{\parskip}{2pt}
\setlength{\emergencystretch}{2em}
\titleformat{\section}{\large\bfseries}{}{0pt}{}[\titlerule]
\titlespacing*{\section}{0pt}{8pt}{4pt}
\setlist[itemize]{leftmargin=14pt,itemsep=2pt,topsep=3pt,parsep=0pt}
\urlstyle{same}
\newcommand{\cvheader}[1]{%
  \begin{minipage}[c]{0.81\textwidth}
  {\LARGE\bfseries Sergio Monge}\\[4pt]
  {\large #1}\\[5pt]
  Trujillo, Peru\\[2pt]
  Phone: \mbox{+51968891526}\\[2pt]
  Email: \mbox{\href{mailto:smonge630@gmail.com}{smonge630@gmail.com}}\\[2pt]
  \href{https://github.com/sergio3212004}{github.com/sergio3212004}\\[2pt]
  \href{https://www.linkedin.com/in/sergio-monge-aa2029306}{linkedin.com/in/sergio-monge-aa2029306}
  \end{minipage}\hfill
  \begin{minipage}[c]{0.16\textwidth}
  \centering
  \IfFileExists{assets/portrait.jpg}{\includegraphics[width=\linewidth,height=3.2cm,keepaspectratio]{assets/portrait.jpg}}{%
    \IfFileExists{assets/portrait.png}{\includegraphics[width=\linewidth,height=3.2cm,keepaspectratio]{assets/portrait.png}}{%
      \IfFileExists{assets/sample-portrait.png}{\includegraphics[width=\linewidth,height=3.2cm,keepaspectratio]{assets/sample-portrait.png}}{}}}
  \end{minipage}\par
}
\newcommand{\translatedsection}[2]{\section{#1}{\footnotesize\itshape #2}\par}
\newcommand{\project}[2]{\textbf{#1}\\{\small #2}\par}

\hypersetup{pdftitle={Sergio Monge - Desarrollador Full Stack},pdfauthor={Sergio Monge},pdflang={es}}
\begin{document}
\cvheader{Desarrollador Full Stack}

\translatedsection{Summary}{Resumen profesional}
Estudiante de Ingeniería Informática en la Universidad Nacional de Trujillo con experiencia en proyectos de desarrollo web full stack. Desarrollo de aplicaciones académicas e institucionales con Laravel, React, TypeScript y PostgreSQL, y colaboración en un editor de recorridos virtuales. Interés en aplicar el desarrollo de software y la resolución de problemas a desafíos prácticos.

\translatedsection{Education}{Formación académica}
\textbf{Universidad Nacional de Trujillo} \hfill Trujillo, Perú\\
Ingeniería Informática --- carrera en curso \hfill 2022--Actualidad

\translatedsection{Projects}{Proyectos}
\project{Sistema de prácticas preprofesionales --- UNT}{PHP, Laravel, PostgreSQL, Vite \enspace |\enspace Proyecto universitario}
\begin{itemize}
\item Desarrollé una aplicación web para gestionar y monitorear las prácticas preprofesionales de la Escuela de Informática de la UNT.
\item Implementé el registro de empresas con consulta de RUC y autocompletado de datos, con ingreso manual cuando el servicio de consulta no está disponible.
\end{itemize}
{\small \href{https://github.com/sergio3212004/sistema-practicas-unt}{github.com/sergio3212004/sistema-practicas-unt}}\par
\vspace{4pt}
\project{Sitio web de Ciencia Política y Gobernabilidad --- UNT}{React, TypeScript, Vite, Tailwind CSS, React Flow}
\begin{itemize}
\item Desarrollé un sitio web institucional con información académica, docentes y noticias para la Escuela de Ciencia Política y Gobernabilidad de la UNT.
\item Incorporé una malla curricular interactiva y una estructura reutilizable para otras escuelas.
\end{itemize}
{\small \href{https://cienciapolitica.unitru.edu.pe}{cienciapolitica.unitru.edu.pe} \enspace |\enspace \href{https://github.com/sergio3212004/pagina-ciencia-politica-unt}{Código fuente en GitHub}}\par
\vspace{4pt}
\project{Tour360 --- Editor colaborativo de recorridos virtuales}{Next.js, React, TypeScript, Three.js, Supabase, Tailwind CSS}
\begin{itemize}
\item Colaboré en el desarrollo de un editor web para crear, editar y publicar recorridos virtuales interactivos de 360 grados.
\end{itemize}
{\small \href{https://esfera360.vercel.app/}{esfera360.vercel.app} \enspace |\enspace \href{https://github.com/NativaMk/tour360-repo}{github.com/NativaMk/tour360-repo}}

\translatedsection{Skills}{Habilidades técnicas}
\textbf{Lenguajes:} PHP, Java, JavaScript, TypeScript, Python, SQL.\\
\textbf{Frontend:} React, Next.js, Vue, Nuxt.js, Angular, Tailwind CSS, Blade.\\
\textbf{Backend:} Laravel, Spring Boot, FastAPI, Django, Express.js, Node.js, NestJS.\\
\textbf{Bases de datos:} PostgreSQL, MySQL, MongoDB.\\
\textbf{Herramientas:} Git, Linux, Docker, Kubernetes.\\
\textbf{Desarrollo asistido por IA:} Codex, Claude, OpenCode.

\translatedsection{Languages}{Idiomas}
\textbf{Inglés:} B2 --- intermedio alto.
\end{document}
